
%%
%% forked from https://gits-15.sys.kth.se/giampi/kthlatex kthlatex-0.2rc4 on 2020-02-13
%% expanded upon by Gerald Q. Maguire Jr.
%% This template has been adapted by Anders Sjögren to the University
%% Engineering Program in Computer Science at KTH ICT. This adaptation was
%% translation of English headings into Swedish with the addition of Swedish.
%% Many thanks to others who have provided constructive input regarding the template.

% Make it possible to conditionally depend on the TeX engine used
\RequirePackage{ifxetex}
\RequirePackage{ifluatex}
\newif\ifxeorlua
\ifxetex\xeorluatrue\fi
\ifluatex\xeorluatrue\fi

\ifxeorlua
% The following is to ensure that the PDF uses a recent version rather than the typical PDF 1-5
%  This same version of PDF should be set as an option for hyperef

\RequirePackage{expl3}
\RequirePackage{etoolbox}


\ExplSyntaxOn
%pdf_version_gset:n{2.0}
%\pdf_version_gset:n{1.5}

%% Alternatively, if you have a LaTeX newer than June 2022, you can use the following. However, then you have to remove the pdfversion from hyperef. It also breaks hyperxmp. So perhaps it is too early to try using it!
%\DocumentMetadata
%{
%% testphase = phase-I, % tagging without paragraph tagging
% testphase = phase-II % tagging with paragraph tagging and other new stuff.
%pdfversion = 2.0 % pdfversion must be set here.
%}

% Optionally, you can set the uncompress flag to make it easier to examine the PDF
%\pdf_uncompress: % to check the pdf
\ExplSyntaxOff
\else
\RequirePackage{expl3}
\ExplSyntaxOn
%\pdf_version_gset:n{2.0}
\pdf_version_gset:n{1.5}
\ExplSyntaxOff
\fi

%% Define a pair of commands to disable and reenable specific packages - see https://tex.stackexchange.com/questions/39415/unload-a-latex-package
\makeatletter
\newcommand{\disablepackage}[2]{%
  \disable@package@load{#1}{#2}%
}
\newcommand{\reenablepackage}[1]{%
  \reenable@package@load{#1}%
}
\makeatother
%% To avoid the warning: "Package transparent Warning: Loading aborted, because pdfTeX is not running in PDF mode."
\ifxeorlua
\disablepackage{transparent}{}
\fi

%% The template is designed to handle a thesis in English or Swedish
\documentclass[english, bibtex]{kththesis}               % the bibliographic processor to use: biblatex or bibtex
%nomenclature,           % add the option nomenclature
%digitaloutput,          % optimize for digital output (this changes the color palette)
%includeExtraAbstracts, % uncomment to allow abstracts in languages beyond English & Swedish


% if pdflatex \usepackage[utf8]{inputenc}

%% Conventions for todo notes:
% Informational
%% \generalExpl{Comments/directions/... in English}
\newcommand*{\generalExpl}[1]{\todo[inline]{#1}}                

% Language-specific information (currently in English or Swedish)
\newcommand*{\engExpl}[1]{\todo[inline, backgroundcolor=kth-lightgreen40]{#1}} %% \engExpl{English descriptions about formatting}
\newcommand*{\sweExpl}[1]{\todo[inline, backgroundcolor=kth-lightblue40]{#1}}  %% % \sweExpl{Text på svenska}

% warnings
\newcommand*{\warningExpl}[1]{\todo[inline, backgroundcolor=kth-lightred40]{#1}} %% \warningExpl{warnings}

% Silence all template-side todo notes so they don't render as colored boxes
% in the final PDF. Uncomment any of these if you want the corresponding notes
% back during editing.
\renewcommand\warningExpl[1]{}
\renewcommand\generalExpl[1]{}
\renewcommand\engExpl[1]{}
\renewcommand\sweExpl[1]{}


% \usepackage[style=numeric,sorting=none,backend=biber]{biblatex}
\iftoggle{biblatex}{
    %\usepackage[language=english,bibstyle=authoryear,citestyle=authoryear, maxbibnames=99]{biblatex}
    % Alternatively you might use another style, such as IEEE and use citestyle=numeric-comp  to put multiple citations in a single pair of square brackets
    \usepackage[style=ieee,citestyle=numeric-comp]{biblatex}
    \addbibresource{references.bib}
    %\DeclareLanguageMapping{norsk}{norwegian}
}{
    % The line(s) below are for BibTeX
    \bibliographystyle{bibstyle/myIEEEtran}
    %\bibliographystyle{apalike}
}


% include a variety of packages that are useful
\input{lib/includes}
\input{lib/kthcolors}

%\glsdisablehyper
%\makeglossaries
%\makenoidxglossaries
%%%% Local Variables:
%%% mode: latex
%%% TeX-master: t
%%% End:
% The following command is used with glossaries-extra
\setabbreviationstyle[acronym]{long-short}
% The form of the entries in this file is \newacronym{label}{acronym}{phrase}
%                                      or \newacronym[options]{label}{acronym}{phrase}

\newacronym{auprc}{AUPRC}{area under the precision-recall curve}
\newacronym{auroc}{AUROC}{area under the receiver operating characteristic curve}
\newacronym{cpu}{CPU}{central processing unit}
\newacronym{json}{JSON}{JavaScript Object Notation}
\newacronym{API}{API}{application programming interface}
\newacronym{bert}{BERT}{Bidirectional Encoder Representations from Transformers}
\newacronym{gpu}{GPU}{graphics processing unit}
\newacronym[longplural={large language models}]{llm}{LLM}{large language model}
\newacronym{llmjudge}{LLM-as-a-Judge}{large language model as a judge}
\newacronym{qa}{QA}{question answering}
\newacronym{rag}{RAG}{Retrieval-Augmented Generation}
\newacronym{tfidf}{TF-IDF}{term frequency-inverse document frequency}
\newacronym{UN}{UN}{United Nations}
\newacronym{SDG}{SDG}{Sustainable Development Goal}
                %load the acronyms file

\input{lib/defines}  % load some additional definitions to make writing more consistent

% The following is needed in conjunction with generating the DiVA data with abstracts and keywords using the scontents package and a modified listings environment
%\usepackage{listings}   %  already included

\makeatletter
\let\verbatimsc\@undefined
\let\endverbatimsc\@undefined
\lst@AddToHook{Init}{\hyphenpenalty=50\relax}
\makeatother


\lstnewenvironment{verbatimsc}
    {
    \lstset{%
        basicstyle=\ttfamily\tiny,
        backgroundcolor=\color{white},
        %basicstyle=\tiny,
        %columns=fullflexible,
        columns=[l]fixed,
        language=[LaTeX]TeX,
        %numbers=left,
        %numberstyle=\tiny\color{gray},
        keywordstyle=\color{red},
        breaklines=true,                 % sets automatic line breaking
        breakatwhitespace=true,          % sets if automatic breaks should only happen at whitespace
        %keepspaces=false,
        breakindent=0em,
        %fancyvrb=true,
        frame=none,                     % turn off any box
        postbreak={}                    % turn off any hook arrow for continuation lines
    }
}{}

%% Add some more keywords to bring out the structure more
\lstdefinestyle{[LaTeX]TeX}{
morekeywords={begin, todo, textbf, textit, texttt}
}

%% definition of new command for bytefield package
\newcommand{\colorbitbox}[3]{%
	\rlap{\bitbox{#2}{\color{#1}\rule{\width}{\height}}}%
	\bitbox{#2}{#3}}




% define a left aligned table cell that is ragged right
\newcolumntype{L}[1]{>{\raggedright\let\newline\\\arraybackslash\hspace{0pt}}p{#1}}

% Because backref is not compatible with biblatex
\iftoggle{biblatex}{
    \usepackage[plainpages=false]{hyperref}
}{
    \usepackage[
    backref=page,
    pagebackref=false,
    plainpages=false,
                            % PDF related options
    unicode=true,           % Unicode encoded PDF strings
    bookmarks=true,         % generate bookmarks in PDF files
    bookmarksopen=false,    % Do not automatically open the bookmarks in the PDF reading program
    pdfpagemode=UseNone,    % None, UseOutlines, UseThumbs, or FullScreen
    destlabel,              % better naming of destinations
    pdfencoding=auto,       % for unicode in 
    ]{hyperref}
   %% Removed use of backref package as it is incompatible with using attachfile
}

\usepackage[all]{hypcap}	%% prevents an issue related to hyperref and caption linking

%% Acronyms
% note that nonumberlist - removes the cross references to the pages where the acronym appears
% note that super will set the descriptions text aligned
% note that nomain - does not produce a main glossary, thus only acronyms will be in the glossary
% note that nopostdot - will prevent there being a period at the end of each entry
\usepackage[acronym, style=super, section=section, nonumberlist, nomain,
nopostdot]{glossaries}
\setlength{\glsdescwidth}{0.75\textwidth}
\usepackage[]{glossaries-extra}
\ifinswedish
    %\usepackage{glossaries-swedish}
\fi

%% For use with the README_notes
% Define a new type of glossary so that the acronyms defined in the README_notes document can be distinct from those in the thesis template
% the tlg, tld, and dn will be the file extensions used for this glossary
\newglossary[tlg]{readme}{tld}{tdn}{README acronyms}


\input{lib/includes-after-hyperref}

%\glsdisablehyper
\makeglossaries
%\makenoidxglossaries

% The following bit of ugliness is because of the problems PDFLaTeX has handling a non-breaking hyphen
% unless it is converted to UTF-8 encoding.
% If you do not use such characters in your acronyms, this could be simplified to just include the acronyms file.
\ifxeorlua
%%% Local Variables:
%%% mode: latex
%%% TeX-master: t
%%% End:
% The following command is used with glossaries-extra
\setabbreviationstyle[acronym]{long-short}
% The form of the entries in this file is \newacronym{label}{acronym}{phrase}
%                                      or \newacronym[options]{label}{acronym}{phrase}

\newacronym{auprc}{AUPRC}{area under the precision-recall curve}
\newacronym{auroc}{AUROC}{area under the receiver operating characteristic curve}
\newacronym{cpu}{CPU}{central processing unit}
\newacronym{json}{JSON}{JavaScript Object Notation}
\newacronym{API}{API}{application programming interface}
\newacronym{bert}{BERT}{Bidirectional Encoder Representations from Transformers}
\newacronym{gpu}{GPU}{graphics processing unit}
\newacronym[longplural={large language models}]{llm}{LLM}{large language model}
\newacronym{llmjudge}{LLM-as-a-Judge}{large language model as a judge}
\newacronym{qa}{QA}{question answering}
\newacronym{rag}{RAG}{Retrieval-Augmented Generation}
\newacronym{tfidf}{TF-IDF}{term frequency-inverse document frequency}
\newacronym{UN}{UN}{United Nations}
\newacronym{SDG}{SDG}{Sustainable Development Goal}
                %load the acronyms file
\else
\input{lib/acronyms-for-pdflatex}
\fi

% place holders for languages lacking .lbx files
\input{lib/placeHolder_lbx_files}

% insert the configuration information with author(s), examiner, supervisor(s), ...
\input{custom_configuration}

% read in the plain text definitions, if they exist
\IfFileExists{custom_configuration_plaintext.tex}{\input{custom_configuration_plaintext.tex}}{}


% Enter the English and Swedish keywords here for use in the PDF metadata _and_ for later use
% following the respective abstract.
% Try to put the words in the same order in both languages to facilitate matching. For example:
\EnglishKeywords{Retrieval-Augmented Generation, Graph RAG, Query routing, LightRAG, 2WikiMultihopQA, LLM-as-a-Judge}
\SwedishKeywords{Retrieval-Augmented Generation, grafbaserad RAG, frågeroutning, LightRAG, 2WikiMultihopQA, LLM-som-domare}

%%%%% For the oral presentation
%% Add this information once your examiner has scheduled your oral presentation
\presentationDateAndTimeISO{}
\presentationLanguage{}  % generally: eng or swe
\presentationRoom{}
\presentationAddress{}
\presentationCity{Stockholm}

% When there are multiple opponents, separate their names with '\&'
% Opponent's information
%\opponentsNames{A. B. Normal \& A. X. E. Normalè}
\opponentsNames{}

% It is possible to have a cover illustration
%\coverIllustration{figures/image1.png}
% if so, the author might acknowledge this or give the copyright information for it with:
%\coverIllustrationCredit{A. B. Normal}

% Once a thesis is approved by the examiner, add the TRITA number
% The TRITA number for a thesis consists of two parts: a series (unique to each school)
% and the number in the series, which is formatted as the year followed by a colon and
% then a unique series number for the thesis - starting with 1 each year.
%\trita{TRITA -- XXX-EX}{2026:0000}

% Put the title, author, and keyword information into the PDF meta information
\input{lib/pdf_related_includes}


% the custom colors and the commands are defined in defines.tex    
\hypersetup{
	colorlinks  = true,
	breaklinks  = true,
	linkcolor   = \linkscolor,
	urlcolor    = \urlscolor,
	citecolor   = \refscolor,
	anchorcolor = black
}

\ifnomenclature
% The following lines make the page numbers and equations hyperlinks in the Nomenclature list
\renewcommand*{\pagedeclaration}[1]{\unskip, \dotfill\hyperlink{page.#1}{page\nobreakspace#1}}
% The following does not work correctly, as the name of the cross-reference is incorrect
%\renewcommand*{\eqdeclaration}[1]{, see equation\nobreakspace(\hyperlink{equation.#1}{#1})}

% You can also change the page heading for the nomenclature
\renewcommand{\nomname}{List of Symbols Used}

% You can even add customization text before the list
\renewcommand{\nompreamble}{The following symbols will be later used within the body of the thesis.}
\makenomenclature
\fi

%
% The commands below are to configure JSON listings
% 
% format for JSON listings
\colorlet{punct}{red!60!black}
\definecolor{delim}{RGB}{20,105,176}
\definecolor{numb}{RGB}{106, 109, 32}
\definecolor{string}{RGB}{0, 0, 0}

\lstdefinelanguage{json}{
    numbers=none,
    numberstyle=\small,
    frame=none,
    rulecolor=\color{black},
    showspaces=false,
    showtabs=false,
    breaklines=true,
    postbreak=\raisebox{0ex}[0ex][0ex]{\ensuremath{\color{gray}\hookrightarrow\space}},
    breakatwhitespace=true,
    basicstyle=\ttfamily\small,
    extendedchars=false,
    upquote=true,
    morestring=[b]",
    stringstyle=\color{string},
    literate=
     *{0}{{{\color{numb}0}}}{1}
      {1}{{{\color{numb}1}}}{1}
      {2}{{{\color{numb}2}}}{1}
      {3}{{{\color{numb}3}}}{1}
      {4}{{{\color{numb}4}}}{1}
      {5}{{{\color{numb}5}}}{1}
      {6}{{{\color{numb}6}}}{1}
      {7}{{{\color{numb}7}}}{1}
      {8}{{{\color{numb}8}}}{1}
      {9}{{{\color{numb}9}}}{1}
      {:}{{{\color{punct}{:}}}}{1}
      {,}{{{\color{punct}{,}}}}{1}
      {\{}{{{\color{delim}{\{}}}}{1}
      {\}}{{{\color{delim}{\}}}}}{1}
      {[}{{{\color{delim}{[}}}}{1}
      {]}{{{\color{delim}{]}}}}{1}
      {’}{{\char13}}1,
}

\lstdefinelanguage{XML}
{
  basicstyle=\ttfamily\color{blue}\bfseries\small,
  morestring=[b]",
  morestring=[s]{>}{<},
  morecomment=[s]{<?}{?>},
  stringstyle=\color{black},
  identifierstyle=\color{blue},
  keywordstyle=\color{cyan},
  breaklines=true,
  postbreak=\raisebox{0ex}[0ex][0ex]{\ensuremath{\color{gray}\hookrightarrow\space}},
  breakatwhitespace=true,
  morekeywords={xmlns,version,type}% list your attributes here
}

% If you use both listing and lstlisting environments - the following makes them both use the same counter and the same formatting in the List of Listings
\makeatletter
\AtBeginDocument{\let\c@listing\c@lstlisting}
\AtBeginDocument{\let\l@listing\l@lstlisting}
\makeatother

% Change the heading for the "List of Listings" to simply "Listings"
\renewcommand{\lstlistlistingname}{Listings}

% number each listing within the chapter
\numberwithin{listing}{chapter}

\usepackage{subfiles}

% To have Creative Commons (CC) license and logos use the doclicense package
% Note that the lowercase version of the license has to be used in the modifier
% i.e., one of by, by-nc, by-nd, by-nc-nd, by-sa, by-nc-sa, zero.
% For background see:
% https://www.kb.se/samverkan-och-utveckling/oppen-tillgang-och-bibsamkonsortiet/open-access-and-bibsam-consortium/open-access/creative-commons-faq-for-researchers.html
% https://kib.ki.se/en/publish-analyse/publish-your-article-open-access/open-licence-your-publication-cc
\begin{comment}
\usepackage[
    type={CC},
    %modifier={by-nc-nd},
    %version={4.0},
    modifier={by-nc},
    imagemodifier={-eu-88x31},  % to get Euro symbol rather than Dollar sign
    hyphenation={RaggedRight},
    version={4.0},
    %modifier={zero},
    %version={1.0},
]{doclicense}
\end{comment}
\listfiles

\begin{document}
%\selectlanguage{swedish}
%
\selectlanguage{english}

%%% Set the numbering for the title page to a numbering series not in the preface or body
\pagenumbering{alph}
\kthcover
\clearpage\thispagestyle{empty}\mbox{} % empty back of front cover
\titlepage

% If you do not want to have a bookinfo page, comment out the line saying \bookinfopage and add a \cleardoublepage
% If you want a bookinfo page: you will get a copyright notice, unless you have used the doclicense package in which case you will get a Creative Commons license. To include the doclicense package, uncomment the configuration of this package above and configure it with your choice of license.
\bookinfopage

% Frontmatter includes the abstracts and table-of-contents
\frontmatter
\setcounter{page}{1}
\begin{abstract}
  \markboth{\abstractname}{}
\begin{scontents}[store-env=lang]
eng
\end{scontents}

\begin{scontents}[store-env=abstracts,print-env=true]
Large language models can generate fluent text but cannot reliably answer questions whose answer lies outside their training data, which has made retrieval-augmented generation a standard architecture for grounded question answering over private corpora. Vector retrieval is the cheap default, but it can fail when the answer requires combining information from several documents. Graph-enhanced retrieval methods such as GraphRAG and LightRAG address this limitation by constructing a graph of entities and relations over the corpus, at the price of being significantly more compute-intensive than vector retrieval. Applying graph-enhanced retrieval to every query is therefore wasteful when the cheap mode already suffices.

This thesis studies whether the choice between LightRAG's cheap vector retrieval mode (\emph{Naive}) and its more expensive graph-enhanced mode (\emph{Mix}) can be predicted from the query text alone, using a lightweight encoder-only classifier as the routing layer rather than a large language-model controller. Routing labels are constructed on the 2WikiMultihopQA benchmark by executing both retrieval modes on each question and using an LLM-as-a-Judge protocol to compare each mode's answer against the gold answer. A ModernBERT-based router is compared against a TF-IDF logistic-regression baseline as well as static always-\emph{Naive}, always-\emph{Mix}, and majority-class reference policies, and is reported with 1\,000-iteration bootstrap confidence intervals across three training seeds.

On the held-out test split, the ModernBERT router reaches an AUPRC of $0.732 \pm 0.005$ on the minority \texttt{mix\_required} class, a balanced accuracy of $0.756 \pm 0.008$, and an F1 of $0.645 \pm 0.002$, improving on a strong TF-IDF baseline (AUPRC 0.708, balanced accuracy 0.727, F1 0.605) while remaining within bootstrap intervals on AUROC. A per-question-type breakdown localizes most of the routing signal in the compositional question slice, where the encoder router improves AUPRC by roughly six points over the TF-IDF baseline. At its validation-selected operating point the router reduces mean routed execution time by roughly one third compared with always-\emph{Mix} while keeping non-trivial recall on the minority class. The result indicates that retrieval-mode selection in a Hybrid retrieval-augmented generation system can be made adaptive without moving the routing decision itself into an expensive model call.
\end{scontents}

\subsection*{Keywords}
\begin{scontents}[store-env=keywords,print-env=true]
Retrieval-Augmented Generation, Graph RAG, Query routing, LightRAG, 2WikiMultihopQA, LLM-as-a-Judge
\end{scontents}
\end{abstract}
\cleardoublepage
\babelpolyLangStart{swedish}
\begin{abstract}
    \markboth{\abstractname}{}
\begin{scontents}[store-env=lang]
swe
\end{scontents}

\begin{scontents}[store-env=abstracts,print-env=true]
Stora språkmodeller kan generera flytande text men kan inte tillförlitligt besvara frågor vars svar ligger utanför deras träningsdata, vilket har gjort hämtningsförstärkt generering (retrieval-augmented generation, RAG) till en standardarkitektur för förankrad frågebesvaring över privata textkorpus. Vektorbaserad hämtning är det billiga standardvalet, men den kan misslyckas när svaret kräver att information från flera dokument kombineras. Grafbaserad RAG, såsom GraphRAG och LightRAG, hanterar detta genom att konstruera en graf över entiteter och relationer i korpusen, till priset av att vara avsevärt mer beräkningsintensiv än vektorhämtning. Att applicera grafbaserad hämtning på varje fråga är därmed slöseri när det billiga läget redan räcker.

Detta examensarbete undersöker om valet mellan LightRAGs billiga vektorläge (\emph{Naive}) och dess dyrare grafförstärkta läge (\emph{Mix}) går att förutsäga enbart från frågetexten, med hjälp av en lättviktig encoder-baserad klassificerare som routinglager istället för en stor språkmodell. Routing-etiketter konstrueras på riktmärket 2WikiMultihopQA genom att köra båda hämtningslägena på varje fråga och använda ett LLM-som-domare-protokoll för att jämföra svaren mot rätt svar. En ModernBERT-baserad router jämförs mot en TF-IDF logistisk regression-baseline samt statiska alltid-\emph{Naive}-, alltid-\emph{Mix}- och majoritetsklasspolicies, med bootstrap-konfidensintervall över tre träningsfrön på testmängden.

På den undanlagda testmängden uppnår ModernBERT-routern AUPRC $0{,}732 \pm 0{,}005$ på minoritetsklassen \texttt{mix\_required}, balanserad noggrannhet $0{,}756 \pm 0{,}008$ och F1 $0{,}645 \pm 0{,}002$, vilket är bättre än en stark TF-IDF-baseline (AUPRC $0{,}708$, balanserad noggrannhet $0{,}727$, F1 $0{,}605$) medan AUROC ligger inom bootstrap-intervallen. En uppdelning per frågetyp visar att större delen av routingsignalen ligger i den kompositionella frågetypen, där encoder-routern förbättrar AUPRC med ungefär sex procentenheter över TF-IDF-baselinen. Vid den validerings\-valda tröskeln minskar routern den genomsnittliga körningstiden med ungefär en tredjedel jämfört med alltid-\emph{Mix} samtidigt som den behåller icke-trivial täckning på minoritetsklassen. Resultatet visar att hämtningslägesval i ett hybrid-RAG-system kan göras adaptivt utan att routningsbeslutet i sig flyttas in i ytterligare en dyr modell.
\end{scontents}

\subsection*{Nyckelord}
\begin{scontents}[store-env=keywords,print-env=true]
Retrieval-Augmented Generation, Grafbaserad RAG, Frågeroutning, LightRAG, 2WikiMultihopQA, LLM-som-domare
\end{scontents}
\end{abstract}
\cleardoublepage
\babelpolyLangStop{swedish}
% Note that the \cleardoublepage has to be before the 
% \babelpolyLangStop - otherwise the running headers will not be in the correct language


\engExpl{If you add the class option \texttt{includeExtraAbstracts} in your \texttt{docuemntclass}, then you can enable abstracts/summaries in a number of languages. Include those versions of abstracts that you need.}
\generalExpl{If you are an exchange student, use the relevant language or languages for abstracts for your home university, as this will often avoid the need for writing another thesis for your home university.}
\generalExpl{If you are fluent in other languages, feel free to add the abstracts in one or more of them.}

\engExpl{The set of languages that can be easily included are based on those languages that were used in theses at KTH in the period 2018-2019, along with several others that have been used in theses since.}
\warningExpl{Note that you may need to augment the set of languages used in \texttt{polyglossia} or
\texttt{babel} (see the file \texttt{kththesis.cls}). If you add a new language, when specifying the language for the abstract, use the three-letter ISO 639-2 Code – specifically the "B" (bibliographic) variant of these codes (note that this is the same language code used in DiVA).}

\cleardoublepage
\iftoggle{includeExtraAbstracts}{
% To include abstracts in languages other than English and Swedish:
%  1. Ensure 'includeExtraAbstracts' is an active option in \documentclass.
%  2. Uncomment the relevant '\input' line(s) below.
%  3. Edit the content within the corresponding file in the 'Additional_Abstracts/' folder.
% Include only the abstracts that you want and fill them in as appropriate
% The template will automatically include the information that you entered
%
% Languages in frequency order based on DiVA data from 2025-05-28

% --- Include Spanish abstract ---
%\input{Additional_Abstracts/abstract_spanish}

% --- Include French abstract ---
%\input{Additional_Abstracts/abstract_french}

% --- Include German abstract ---
%\input{Additional_Abstracts/abstract_german}

% --- Include Italian abstract ---
%\input{Additional_Abstracts/abstract_italian}

% --- Include Chinese abstract ---
%\input{Additional_Abstracts/abstract_chinese}

% --- Include Norwegian abstract ---
%\input{Additional_Abstracts/abstract_norweign}

% --- Include Catalan abstract ---
%\input{Additional_Abstracts/abstract_catalan}

% --- Include Finnish abstract ---
%\input{Additional_Abstracts/abstract_finnish}

% --- Include Portuguese abstract ---
%\input{Additional_Abstracts/abstract_chinese}

% --- Include Dutch abstract ---
%\input{Additional_Abstracts/abstract_dutch}

% --- Include Swahili / Kiswahili summary ---
%\input{Additional_Abstracts/abstract_swahili}

% --- Include Russian summary ---
%\input{Additional_Abstracts/abstract_russian}

% --- Include Polish abstract ---
%\input{Additional_Abstracts/abstract_polishn}

% --- Include Ukrainian abstract ---
%\input{Additional_Abstracts/abstract_ukrainian}

% --- Include Hindi abstract ---
%\input{Additional_Abstracts/abstract_hindi}

% --- Include Danish abstract ---
%\input{Additional_Abstracts/abstract_danish}

% --- Include Estonian abstract ---
%\input{Additional_Abstracts/abstract_estonian}

% --- Include Slovak abstract ---
%\input{Additional_Abstracts/abstract_slovak}

% --- Include Serbian abstract ---
%\input{Additional_Abstracts/abstract_serbian}

% --- Include Bosnian abstract ---
%\input{Additional_Abstracts/abstract_bosnian}

%% missing many languages



}{} %% End the iiftoggle{includeExtraAbstracts}

\section*{Acknowledgments}
\markboth{Acknowledgments}{}

I would like to thank my supervisor at KTH, Arvid Eriksson, for his detailed and patient feedback on the pre-study and the report drafts, and for steering the project toward the right level of scientific rigor at every step. I would also like to thank my external supervisor at Ericsson, Peng Zhang, for the industrial context that motivated the routing problem, and my examiner at KTH, Pawel Herman, for accepting the project and shaping its overall framing.

Generative AI was used as a writing partner during structural revision and code-level analysis of the experimental results. All scientific decisions, technical claims, experimental code, and final formulations have been manually reviewed and accepted by the author.

\acknowlegmentssignature

\fancypagestyle{plain}{}
\renewcommand{\chaptermark}[1]{ \markboth{#1}{}} 
\tableofcontents
  \markboth{\contentsname}{}

\cleardoublepage
\listoffigures

\cleardoublepage

\listoftables
\cleardoublepage
\lstlistoflistings\engExpl{If you have listings in your thesis. If not, then remove this preface page.}
\cleardoublepage
% Align the text expansion of the glossary entries
\newglossarystyle{mylong}{%
  \setglossarystyle{long}%
  \renewenvironment{theglossary}%
     {\begin{longtable}[l]{@{}p{\dimexpr 2cm-\tabcolsep}p{0.8\hsize}}}% <-- change the value here
     {\end{longtable}}%
 }
%\glsaddall
%\printglossaries[type=\acronymtype, title={List of acronyms}]
\printglossary[style=mylong, type=\acronymtype, title={List of acronyms and abbreviations}]
%\printglossary[type=\acronymtype, title={List of acronyms and abbreviations}]

%\printnoidxglossary[style=mylong, title={List of acronyms and abbreviations}]
\engExpl{The list of acronyms and abbreviations should be in alphabetical order based on the spelling of the acronym or abbreviation.
}

% if the nomenclature option was specified, then include the nomenclature page(s)
\ifnomenclature
    \cleardoublepage
    % Output the nomenclature list
    \printnomenclature
\fi

%% The following label is essential to know the page number of the last page of the preface
%% It is used to compute the data for the "For DIVA" pages
\label{pg:lastPageofPreface}
% Mainmatter is where the actual contents of the thesis goes
\mainmatter
\glsresetall
\renewcommand{\chaptermark}[1]{\markboth{#1}{}}
\selectlanguage{english}
\chapter{Introduction}
\label{ch:introduction}

This chapter introduces the thesis topic and motivates the problem studied in this thesis project. Section~\ref{sec:background} gives the necessary background in retrieval-augmented question answering and the cost of different retrieval modes. Section~\ref{sec:problem} defines the problem and research question. Section~\ref{sec:purpose} states the purpose of the work, Section~\ref{sec:goals} presents the goals, Section~\ref{sec:researchMethodology} outlines the research methodology, Section~\ref{sec:delimitations} defines the delimitations, and Section~\ref{sec:structure} describes the structure of the rest of the thesis.

\section{Background}
\label{sec:background}

Large language models can produce fluent answers but they cannot reliably answer questions whose answer lies outside their training data, and they often hallucinate when asked about private or recent information \cite{gao2023ragsurvey}. \Gls{rag} addresses this by retrieving relevant passages from an external corpus and conditioning the language model on those passages, which lets the model cite specific sources and reduces hallucination on private or domain-specific questions \cite{gao2023ragsurvey,lewis2020rag}. This makes \gls{rag} the dominant architecture for grounded question answering on company-internal documents, which is also the industrial setting that motivates this thesis at Ericsson.

The simplest retrieval method is vector similarity search. Documents are split into chunks, each chunk is embedded into a vector space using an embedding model, and the chunks closest to the query embedding are retrieved. This works well when the answer is contained inside a single passage that resembles the query, but it can fail when the answer requires combining facts from multiple documents. Multi-hop question answering benchmarks are designed around exactly this regime, in which the answer can only be derived by chaining several supporting facts together. 2WikiMultihopQA \cite{ho2020twowiki} is one such benchmark and is the dataset used in this thesis. Graph-enhanced retrieval methods such as GraphRAG and LightRAG address this multi-hop weakness by additionally constructing a graph of entities and relations over the corpus, so that retrieval can follow the structure of the data rather than only local passage similarity \cite{edge2024graphrag,lightrag2024}.

Graph-enhanced retrieval is, however, significantly more computationally expensive than vector retrieval. Building the graph itself requires many language-model calls for entity and relation extraction, and answering a query involves traversing the graph and assembling a larger context. Microsoft's own cost discussion identifies these repeated language-model calls as the dominant cost driver in GraphRAG \cite{edge2024graphrag,microsoft2024graphragcosts}. Graph-enhanced retrieval is therefore best understood as an expensive capability that should be used when it is actually needed.

This motivates adaptive retrieval. Rather than applying the same retrieval strategy to every query, the system should choose a strategy based on the query. Self-RAG learns when to retrieve through self-reflection, and Adaptive-RAG selects a retrieval strategy based on an estimate of question complexity \cite{asai2023selfrag,jeong2024adaptiverag}. More generally, model routing decides at the question level whether to send the request to a powerful but expensive language model or to a smaller, cheaper one with weaker performance, and routing benchmarks such as RouterBench frame this as a cost--quality trade-off \cite{hu2024routerbench}. This thesis adopts the same framing but applies it inside the retrieval layer rather than across language models, deciding whether each query should be served by a cheap vector path or by a more expensive graph-enhanced path.

\section{Problem}
\label{sec:problem}

A static retrieval policy in a Hybrid \gls{rag} system is suboptimal. Always using the cheap vector path can miss queries whose answer requires structured evidence from several documents, and always using the expensive graph-enhanced path wastes computation on queries that the cheap path already answers correctly. The research problem is to decide, for a given query and before retrieval is executed, which retrieval path should be used.

The decision is made harder by the fact that the routing decision should itself be cheap. If a large general-purpose language model is used as the controller, much of the intended cost saving is lost at the control layer. Prior work on encoder-only text classification shows that compact bidirectional models can match larger decoder models on narrow classification tasks while remaining significantly cheaper to run, which makes them attractive candidates for the routing layer \cite{benayas2024encoder}. The open question is whether the vector-sufficient versus graph-required distinction is predictable from the query text alone, well enough for such a lightweight classifier to be useful.

\subsection{Research question}
\label{sec:researchQuestion}

Within a fixed Hybrid \gls{rag} setting, how well can a lightweight query-only classifier predict whether a query needs graph-enhanced retrieval rather than cheap vector retrieval, measured by held-out \gls{auprc} and by the cost--performance trade-off against always-vector and always-graph baselines?

\section{Purpose}
\label{sec:purpose}

The purpose of this thesis is to investigate whether retrieval-mode selection in a Hybrid \gls{rag} system can be made adaptive without moving the routing decision itself into an expensive model call, so that graph-enhanced retrieval is used when it is needed and avoided when vector retrieval already suffices. The contribution is empirical. It is a quantitative account, on a fixed public benchmark, of how predictable retrieval need is from query text alone and how that predictability translates into routed cost.

\section{Goals}
\label{sec:goals}

The research question is broken down into four concrete goals.

\begin{itemize}
    \item Build a reproducible pipeline that prepares 2WikiMultihopQA samples, indexes them in LightRAG, and executes both LightRAG's \emph{Naive} and \emph{Mix} retrieval modes per query while logging execution time.
    \item Construct a supervised routing dataset by labeling each query with the minimum viable retrieval mode, using \gls{llmjudge} to compare each retrieval mode's answer against the benchmark answer.
    \item Train a lightweight query-only router and compare it against a classical \gls{tfidf} logistic-regression baseline as well as static always-\emph{Naive} and always-\emph{Mix} policies.
    \item Evaluate the router both as a held-out classifier, using metrics such as \gls{auprc}, precision, recall, and F1, and as a thresholded routing policy, using average routed execution time against the static policies.
\end{itemize}

\section{Research methodology}
\label{sec:researchMethodology}

This thesis is a quantitative empirical study on a public benchmark, using held-out classification metrics and offline policy simulation. The 2WikiMultihopQA corpus is indexed in LightRAG, every benchmark question is executed in both retrieval modes, and \gls{llmjudge} assigns each query a routing label based on whether the cheap mode was already sufficient. Bootstrap confidence intervals are reported on the held-out test set so that differences between the router and the baselines can be compared rather than read off point estimates.

\section{Delimitations}
\label{sec:delimitations}

The router is query-only and decides retrieval mode before any retrieval is run. It does not inspect retrieved passages, intermediate graph state, or any draft answer. The retrieval framework is LightRAG and the two routed modes are LightRAG's \emph{Naive} mode, used as the cheap vector path, and its \emph{Mix} mode, used as the graph-enhanced path. Results about LightRAG's mode selection should not be assumed to transfer automatically to other Hybrid \gls{rag} frameworks.

The benchmark is 2WikiMultihopQA \cite{ho2020twowiki}, chosen because it is public and provides supporting evidence, which makes the experiment reproducible. Ericsson documentation motivates the industrial relevance of the work but is not used as a benchmark, since the unstructured corpus lacks the per-query ground truth needed for judged routing labels.

Routing labels are produced by \gls{llmjudge} and are treated as judged supervision rather than absolute ground truth. Queries where neither retrieval mode produces a correct answer are excluded from router training.

\cleardoublepage
\chapter{Background}
\label{ch:background}

This chapter introduces the technical and scientific concepts that the rest of the thesis builds on. Section~\ref{sec:backgroundRag} describes \gls{rag}, why it is needed for grounding large language model outputs, and how the components fit together. Section~\ref{sec:backgroundHybridGraph} contrasts vector retrieval with graph-enhanced retrieval, including the two LightRAG retrieval modes used in this thesis. Section~\ref{sec:backgroundAdaptiveRetrieval} describes adaptive retrieval and query routing. Section~\ref{sec:backgroundEncoderClassifiers} introduces the language models and text classifiers used by the router. Section~\ref{sec:relatedWork} positions the thesis in relation to prior work.

\section{Retrieval-augmented generation}
\label{sec:backgroundRag}

Large language models encode an enormous amount of general knowledge in their parameters, but they have two limitations relevant to industrial question answering. They cannot access information that was not part of their training data, which makes them unable to answer questions over private or company-internal documents, and they have a tendency to hallucinate, generating fluent but factually incorrect answers when asked about content outside their training distribution \cite{gao2023ragsurvey}. \Gls{rag} addresses both problems by retrieving relevant passages from an external corpus at query time and conditioning the language model on those passages, which both gives the model access to information beyond its training data and lets it ground its answers in cite-able sources \cite{lewis2020rag,gao2023ragsurvey}.

The original \gls{rag} formulation by Lewis et al.\ \cite{lewis2020rag} treats retrieval and generation jointly as a latent-variable problem. Given a query \(q\), a retriever returns a set of documents \(\mathcal{D}_k(q)\), and the generator produces an answer \(a\) by marginalizing over which document was retrieved, as in
\begin{equation}
\label{eq:ragGeneration}
p(a \mid q) =
\sum_{d \in \mathcal{D}_k(q)}
p_\theta(a \mid q,d)\,p_\eta(d \mid q).
\end{equation}
Here \(p_\eta(d \mid q)\) is the retriever, parameterized by \(\eta\), which assigns a probability to each candidate document given the query, and \(p_\theta(a \mid q, d)\) is the generator, parameterized by \(\theta\), which produces the answer given the query and a retrieved document. Equation~\ref{eq:ragGeneration} is not how every modern \gls{rag} system is implemented in practice, but it makes one point that matters for this thesis explicit. Which documents the retrieval system surfaces directly determines what the generator can answer. Retrieval is part of the system's behavior, not a preprocessing step.

For industrial question answering the motivation is concrete. At Ericsson, internal technical documents cannot be assumed to appear in any public language-model pretraining corpus, so a \gls{rag} system is needed to ground answers in company-controlled knowledge. The same is true of many other domain-specific settings, which is why \gls{rag} has become a standard architecture for grounded question answering over private corpora \cite{gao2023ragsurvey}.

Many real questions cannot be answered from a single passage. Multi-hop question answering benchmarks such as HotpotQA, 2WikiMultihopQA, and MuSiQue are constructed around questions whose answers require combining information from multiple supporting documents \cite{yang2018hotpotqa,ho2020twowiki,trivedi2022musique}. For example, a question of the form \emph{``Who is the spouse of the director of film X?''} requires the system first to identify the director of \(X\), then to look up that person, and finally to retrieve their spouse. A retriever that scores each document independently against the query may surface only part of the chain.

\section{Graph-enhanced retrieval}
\label{sec:backgroundHybridGraph}

Most practical \gls{rag} systems start with vector retrieval. Documents are split into chunks, each chunk is encoded into a dense vector by an embedding model, and at query time the chunks whose vectors are closest in the embedding space to the query's vector are returned to the generator. The original \gls{rag} model used dense vector retrieval over Wikipedia passages \cite{lewis2020rag}, and dense passage retrieval generalized this approach to open-domain question answering \cite{karpukhin2020dpr}. Vector retrieval is fast and works well when the answer is contained inside a single passage that resembles the query.

Vector retrieval is weaker when the required evidence is split across multiple documents or depends on relationships between entities. Closeness in the embedding space measures local textual similarity, not the structure that links documents together. As a result, a nearest-neighbor lookup may retrieve passages that are individually relevant but collectively insufficient to answer a multi-hop question.

\subsection{Graph-based RAG systems}
\label{sec:backgroundGraphSystems}

Graph-based \gls{rag} systems address this limitation by constructing a graph over the corpus before any query is asked. Nodes correspond to entities, chunks, or higher-level summaries, and edges correspond to relations between them. At query time the system can then retrieve evidence by traversing this graph rather than only matching the query against individual chunks.

GraphRAG \cite{edge2024graphrag} extracts entities and their relations from each chunk using a language model, builds a knowledge graph over the entire corpus, and groups closely related entities into communities. Each community is summarized into a higher-level description, which can be used to answer queries that span large parts of the corpus. LightRAG \cite{lightrag2024} follows a similar two-level idea, but is designed to be faster and easier to deploy. It also extracts an entity-relation graph during indexing and uses it together with vector chunks at retrieval time, and it exposes several retrieval modes that vary in how heavily they use this graph. The key practical difference for this thesis is that LightRAG explicitly exposes a cheap vector-only retrieval mode and a more expensive graph-enhanced mode in the same framework, which makes it possible to compare the two modes on identical inputs.

Graph-enhanced retrieval is significantly more computationally expensive than pure vector retrieval. Building the graph requires many language-model calls during indexing for entity and relation extraction, and answering a query involves traversing the graph, assembling a larger context, and often additional language-model calls. Microsoft's cost discussion for GraphRAG identifies these repeated language-model calls as the dominant cost driver \cite{edge2024graphrag,microsoft2024graphragcosts}. Graph-enhanced retrieval is therefore best treated as an expensive capability that is worth invoking when the question benefits from it, and worth avoiding otherwise.

\subsection{LightRAG \emph{Naive} and \emph{Mix} modes}
\label{sec:backgroundLightragModes}

LightRAG exposes several retrieval modes, two of which are used as the routed paths in this thesis \cite{lightrag2024}. The \emph{Naive} mode retrieves the top-\(k\) text chunks for the query by vector similarity and passes them directly to the generator, ignoring the graph. The \emph{Mix} mode is hybrid. It retrieves vector chunks as in \emph{Naive} but also retrieves entities and relations from the entity-relation graph, then assembles all of this into the generator's context. The two modes therefore share the same underlying corpus and generator but differ in how much graph structure they exploit, which makes them a natural pair of routed retrieval paths.

Figure~\ref{fig:vectorVsGraphExample} illustrates the difference with a stylized multi-hop example. The \emph{Naive} mode retrieves chunks similar to the query and can fail when no single chunk contains both required hops. The \emph{Mix} mode follows entity-relation edges in the graph and can connect the two hops even when they live in separate documents.

\begin{figure}[!ht]
\centering
\resizebox{0.95\textwidth}{!}{%
\begin{tikzpicture}[
    every node/.style={font=\small},
    chunk/.style={draw, rounded corners, align=center, minimum width=2.6cm, minimum height=0.85cm, fill=gray!10},
    entity/.style={draw, circle, align=center, minimum size=0.85cm, fill=kth-lightblue40},
    edge/.style={->, thick},
    weakedge/.style={->, thick, dashed, gray}
]
    \node[align=center] (qa) at (0,2.0) {\textbf{Naive (vector only)}};
    \node[align=center] (qq) at (0,1.3) {Query:\\\emph{Spouse of director of film X?}};
    \node[chunk] (c1) at (-1.5,0) {Chunk: film X\\(plot summary)};
    \node[chunk] (c2) at (1.5,0) {Chunk: film X\\(release info)};
    \node[align=center, text width=4cm] (naiveAns) at (0,-1.4) {Retrieved chunks mention\\film X but not its director\\$\Rightarrow$ answer incomplete};
    \draw[weakedge] (qq) -- (c1);
    \draw[weakedge] (qq) -- (c2);
    \draw[weakedge] (c1) -- (naiveAns);
    \draw[weakedge] (c2) -- (naiveAns);

    \node[align=center] (qb) at (8,2.0) {\textbf{Mix (graph-enhanced)}};
    \node[align=center] (qq2) at (8,1.3) {Query:\\\emph{Spouse of director of film X?}};
    \node[entity] (filmX) at (6.4,0) {film X};
    \node[entity] (dir) at (8,0) {director};
    \node[entity] (sp) at (9.6,0) {spouse};
    \node[align=center, text width=4cm] (mixAns) at (8,-1.4) {Two-hop walk on\\entity-relation graph\\$\Rightarrow$ correct answer};
    \draw[edge] (filmX) -- node[above,font=\scriptsize]{directed-by} (dir);
    \draw[edge] (dir) -- node[above,font=\scriptsize]{married-to} (sp);
    \draw[weakedge] (qq2) -- (filmX);
    \draw[weakedge] (sp) -- (mixAns);
\end{tikzpicture}%
}
\caption{Stylized illustration of vector retrieval (\emph{Naive}) versus graph-enhanced retrieval (\emph{Mix}) on a two-hop question. Vector retrieval matches chunks individually against the query, while the graph-enhanced mode can follow entity-relation edges that connect facts living in different documents.}
\label{fig:vectorVsGraphExample}
\end{figure}

\section{Adaptive retrieval and query routing}
\label{sec:backgroundAdaptiveRetrieval}

Adaptive retrieval rests on a simple observation. Not every query benefits from the same retrieval strategy. Self-RAG \cite{asai2023selfrag} fine-tunes the generator itself to emit reflection tokens that decide whether retrieval is necessary, what to retrieve, and whether the generated answer is supported by the evidence, which lets the system skip retrieval entirely on questions it is confident about. Adaptive-RAG \cite{jeong2024adaptiverag} takes a slightly different approach. It trains a small classifier to estimate the complexity of a question and uses that estimate to select between progressively more expensive retrieval-augmented reasoning strategies, ranging from no retrieval through single-hop retrieval to iterative retrieval. Both works support the premise that retrieval behavior should depend on the query, and Adaptive-RAG in particular shows that a small upstream classifier can drive the strategy choice.

Model routing frames the same idea more generally as a cost--quality decision \cite{hu2024routerbench}. In a multi-model setting, the system has access to several language models with different cost and accuracy profiles, and a router decides on a per-question basis whether to use a powerful but expensive model or a cheaper, weaker one. RouterBench \cite{hu2024routerbench} is a benchmark and evaluation suite for this problem and reports that even simple routers can recover a large fraction of the strongest model's accuracy while spending only a fraction of its cost. Although RouterBench concerns routing between language models rather than between retrieval modes, the abstraction transfers directly. A router observes an input and chooses among execution paths with different cost and expected performance before the expensive computation happens.

This thesis applies that routing abstraction inside the retrieval layer. The action space is restricted to two retrieval modes, LightRAG's \emph{Naive} and \emph{Mix}, inside one fixed Hybrid \gls{rag} framework. Figure~\ref{fig:retrievalModeRouting} illustrates the routing decision. The router observes only the query and selects either the cheap vector path or the graph-enhanced path before retrieval runs. Restricting the router to the query is a deliberate design choice. A router that has to run retrieval first in order to decide would no longer be a cheap control layer, and the potential efficiency gain would be lost. Adaptive-RAG and similar work also make their routing decision before retrieval, on the query alone \cite{jeong2024adaptiverag,hu2024routerbench}.

\begin{figure}[!ht]
\centering
\resizebox{\textwidth}{!}{%
\begin{tikzpicture}[
    every node/.style={font=\small},
    box/.style={draw, rounded corners, align=center, minimum width=2.8cm, minimum height=0.95cm},
    pathbox/.style={draw, rounded corners, align=center, minimum width=3.6cm, minimum height=1.15cm},
    decision/.style={draw, diamond, aspect=1.8, align=center, inner sep=2pt, minimum width=2.2cm},
    line/.style={->, thick}
]
    \node[box] (query) at (0,0) {User query\\\(q\)};
    \node[decision] (router) at (3.1,0) {Query\\router};
    \node[pathbox] (naive) at (6.6,1.35) {\textbf{Vector sufficient}\\Low-cost path\\LightRAG \emph{Naive}};
    \node[pathbox] (mix) at (6.6,-1.35) {\textbf{Graph required}\\High-cost path\\LightRAG \emph{Mix}};
    \node[box] (answer) at (10.5,0) {Generated\\answer};

    \draw[line] (query) -- (router);
    \draw[line] (router) -- (naive.west);
    \draw[line] (router) -- (mix.west);
    \draw[line] (naive) -- (answer);
    \draw[line] (mix) -- (answer);
\end{tikzpicture}%
}
\caption{The retrieval-mode routing problem studied in this thesis. The router observes only the query and selects between LightRAG's cheap vector path (\emph{Naive}) and its graph-enhanced path (\emph{Mix}) before any retrieval is run.}
\label{fig:retrievalModeRouting}
\end{figure}

\section{Language models and text classifiers}
\label{sec:backgroundEncoderClassifiers}

The routing layer is implemented as a text classifier over the query. This section describes the two families of classifier compared in this thesis. The first is a transformer-based encoder model, ModernBERT. The second is a classical \gls{tfidf} logistic-regression baseline.

\subsection{Encoder-only transformers and ModernBERT}
\label{sec:backgroundEncoders}

The transformer architecture \cite{devlin2019bert} introduced self-attention as a general mechanism for contextualized text representations. Encoder-only models such as \gls{bert} are trained on masked language modeling and produce a contextualized vector per input token, which can be pooled or fed to a small classification head to produce a class label \cite{devlin2019bert}. Decoder-only models such as the GPT family are instead trained to predict the next token and are more naturally used as generative language models.

Both families can in principle be used for classification, but encoder-only models remain attractive when the output is a class label and inference efficiency matters. A recent comparison by Benayas, Sicilia, and Mora-Cantallops \cite{benayas2024encoder} finds that encoder-only models match or exceed much larger decoder-only models on intent classification and sentiment analysis while being orders of magnitude smaller, which is consistent with the design goal of a cheap control layer.

ModernBERT \cite{warner2024modernbert} is a recent encoder model that updates the \gls{bert} architecture with modern training data, longer context, and several efficiency improvements. The authors report that ModernBERT-base matches or exceeds the classification quality of much larger models while remaining cheap enough to fine-tune on a single \gls{gpu}, which makes it a natural choice for a query-only router that has to remain inexpensive relative to the retrieval modes it routes between.

\subsection{TF-IDF logistic regression}
\label{sec:backgroundTfidf}

The classical baseline used in this thesis is a logistic regression over \gls{tfidf} features. Each query is represented by a sparse vector whose entries are term frequencies weighted by the inverse document frequency of the term in the training corpus, which downweights very common words. A logistic regression then maps this sparse feature vector to a probability for the \texttt{mix\_required} class. The baseline matters for the routing problem because the distinction between vector-sufficient and graph-required queries may be partly lexical. Certain question forms, such as \emph{``the spouse of the director of \ldots''}, \emph{``who first \ldots''}, or comparisons between entities, may simply correlate with graph need. If a bag-of-words classifier can recover most of the signal, then a transformer encoder buys little.

\section{Related work}
\label{sec:relatedWork}

The most closely related prior work falls into three groups. The first is graph-based \gls{rag} systems. The second is adaptive-retrieval work that uses an upstream classifier. The third is general model routing.

\emph{Graph-based RAG.} GraphRAG \cite{edge2024graphrag} and LightRAG \cite{lightrag2024} demonstrate that graph-enhanced retrieval can add structure beyond local vector similarity and improve answer quality on questions that require linking facts across documents. They also document the additional cost of doing so. Neither work, however, asks the next question. \emph{Given} that graph-enhanced retrieval exists alongside cheaper vector retrieval, how should the system decide which one to use for a given query? The work in this thesis treats that selection as a supervised prediction problem.

\emph{Adaptive retrieval with an upstream classifier.} Adaptive-RAG \cite{jeong2024adaptiverag} is the closest precedent on the methodological side. It trains a small classifier on the question to estimate complexity and uses the estimate to select between progressively more expensive retrieval-augmented reasoning strategies, ranging from no retrieval to iterative retrieval. The classifier is trained on labels derived from how well different strategies answer the question on a training set, which is conceptually similar to the judge-derived labels used here. The differences are that Adaptive-RAG selects across reasoning strategies at the language-model level (no retrieval, single retrieval, iterative retrieval), while this thesis selects between two retrieval modes inside a fixed framework (\emph{Naive} versus \emph{Mix}), and that the classifier here is intentionally restricted to a lightweight encoder so the routing cost remains negligible. Self-RAG \cite{asai2023selfrag} achieves a similar adaptive behavior by training the generator itself to emit reflection tokens, which is a more invasive change and does not give an explicit pre-retrieval routing decision.

\emph{Model routing.} RouterBench \cite{hu2024routerbench} formalizes routing between language models as a cost--quality problem and provides a benchmark for it. The setting is different (different models rather than different retrieval modes), but the framing carries over directly. The cost--performance trade-off plots used in Chapter~\ref{ch:implementation} follow the same logic. Each operating point of the learned router can be compared to static baselines, and the relevant question is whether the router lies above the trivial straight line between always-cheap and always-expensive policies.

What is missing from this literature is a focused empirical study of the \emph{Naive}-versus-\emph{Mix} retrieval-mode choice inside a fixed Hybrid \gls{rag} framework as a supervised, query-only routing problem with judge-derived labels. That is the gap this thesis fills.

\cleardoublepage
\chapter{Methods}
\label{ch:methods}

This chapter describes the data, the routing pipeline, the router models, and the evaluation procedure used to study query-only retrieval-mode routing. Section~\ref{sec:methodData} describes the benchmark data and how it is preprocessed and indexed. Section~\ref{sec:methodPipeline} describes the routing-label pipeline that turns raw benchmark questions into a supervised routing dataset. Section~\ref{sec:methodDataset} describes the resulting binary routing dataset and the held-out splits used for evaluation. Section~\ref{sec:methodRouters} describes the two router models, a ModernBERT-based classifier and a \gls{tfidf} logistic-regression baseline. Section~\ref{sec:methodEvaluation} describes the held-out classification metrics and the offline routing-policy evaluation.

\section{Data}
\label{sec:methodData}

\subsection{2WikiMultihopQA}
\label{sec:method2wiki}

The benchmark used in this thesis is 2WikiMultihopQA \cite{ho2020twowiki}, a public multi-hop question answering dataset constructed from Wikipedia. Each sample contains a question, a gold answer, a set of context documents drawn from Wikipedia, and a set of supporting facts that mark which sentences in the context are required to derive the answer. Questions are annotated with one of four types that describe the reasoning structure. These types are \texttt{comparison}, \texttt{compositional}, \texttt{bridge\_comparison}, and \texttt{inference}.

2WikiMultihopQA is chosen for three reasons. It is public, which makes the experiment reproducible. It is multi-hop by construction, so the answer cannot in general be read off a single passage, which is the setting in which graph-enhanced retrieval is expected to differ most from vector retrieval. And the supporting-fact annotations make it possible to relate the routing decision back to which documents the answer actually depends on. Other multi-hop benchmarks such as HotpotQA and MuSiQue \cite{yang2018hotpotqa,trivedi2022musique} share this motivation. 2WikiMultihopQA was preferred because its context documents are short and well-suited to indexing in LightRAG without further chunking heuristics.

\subsection{Preprocessing and indexing}
\label{sec:methodPreprocessing}

The benchmark splits are processed in their original sample order. Each sample is converted into a structured record containing a stable sample identifier, the question, the gold answer, the question type, the context documents, and the supporting facts. The mapping between sample identifier and these fields is preserved throughout the pipeline so that every later artifact (retrieval output, judge label, router prediction) can be traced back to the original sample.

Context documents from these structured records are then inserted into LightRAG \cite{lightrag2024}, which during indexing builds both a vector index over chunked text and an entity-relation graph derived from the same chunks. The indexing pass is incremental and idempotent. Previously inserted documents are recognized and skipped, which makes it possible to resume an interrupted run without rebuilding the index from scratch. Once the index has been built, it is treated as a fixed part of the experimental setup. Stochasticity in graph construction (entity extraction, relation extraction, summarization) is therefore baked into the index and does not introduce additional variance during the dual-mode retrieval experiments.

\section{Routing-label pipeline}
\label{sec:methodPipeline}

Figure~\ref{fig:methodPipeline} shows the end-to-end pipeline. Starting from 2WikiMultihopQA samples, the corpus is indexed in LightRAG, every queryable sample is executed with both \emph{Naive} and \emph{Mix} retrieval, an \gls{llmjudge} assigns a routing class to each query, and the resulting binary labels are used to train and evaluate query-only routers.

\begin{figure}[!ht]
\centering
\resizebox{\textwidth}{!}{%
\begin{tikzpicture}[
    node distance=1.0cm and 1.0cm,
    every node/.style={font=\small},
    box/.style={draw, rounded corners, align=center, minimum width=3.0cm, minimum height=1.0cm},
    line/.style={->, thick}
]
    \node[box] (data) {2WikiMultihopQA\\samples};
    \node[box, right=of data] (index) {LightRAG\\index};
    \node[box, right=of index] (query) {Run \emph{Naive}\\and \emph{Mix}};
    \node[box, right=of query] (judge) {LLM-as-a-Judge\\routing label};
    \node[box, below=of judge] (dataset) {Binary router\\dataset};
    \node[box, left=of dataset] (train) {Train query-only\\routers};
    \node[box, left=of train] (eval) {Held-out and\\policy evaluation};

    \draw[line] (data) -- (index);
    \draw[line] (index) -- (query);
    \draw[line] (query) -- (judge);
    \draw[line] (judge) -- (dataset);
    \draw[line] (dataset) -- (train);
    \draw[line] (train) -- (eval);
\end{tikzpicture}%
}
\caption{End-to-end pipeline from 2WikiMultihopQA samples to router evaluation. Each step writes its output to disk so the remaining steps can be re-run without repeating expensive earlier steps.}
\label{fig:methodPipeline}
\end{figure}

\subsection{Dual-mode retrieval}
\label{sec:methodDualMode}

For every benchmark question \(q_i\), the indexed corpus is queried with both LightRAG retrieval modes, giving
\begin{equation}
\label{eq:dualModeOutputs}
r_i^{N} = R_N(q_i),
\qquad
r_i^{M} = R_M(q_i),
\end{equation}
where \(R_N\) denotes LightRAG \emph{Naive} (the cheap vector path) and \(R_M\) denotes LightRAG \emph{Mix} (the graph-enhanced path). The recorded artifacts for each query are the two generated answers \(r_i^N\) and \(r_i^M\), the measured execution time \(t_i^N\) and \(t_i^M\) for each mode, and the number of context documents and supporting facts associated with the sample.

Execution time is measured per query and used as the cost metric throughout the thesis. Token usage would be a natural alternative cost dimension, but LightRAG's modes invoke an internal language model several times during retrieval (for example for graph traversal and context assembly) and these internal calls were not instrumented in this work. The choice of execution time as the cost metric is discussed further in Section~\ref{sec:methodEvaluation} and revisited in Chapter~\ref{ch:conclusionsAndFutureWork}.

\subsection{LLM-as-a-Judge labeling}
\label{sec:methodJudge}

The 2WikiMultihopQA benchmark provides gold answers but does not provide routing labels. Routing labels are constructed by running an \gls{llmjudge} \cite{zheng2023judge} that compares each query's two generated answers against the benchmark answer. For each query the judge returns one of three classes. The class \texttt{naive\_enough} is assigned when \(r_i^N\) matches the gold answer, so the cheap mode was sufficient. The class \texttt{mix\_required} is assigned when \(r_i^N\) does not match the gold answer but \(r_i^M\) does, so the graph-enhanced mode was required to recover the answer. The class \texttt{none\_enough} is assigned when neither mode produced an answer that matches the gold answer.
This produces the binary routing label
\begin{equation}
\label{eq:routingLabel}
y_i =
\begin{cases}
0, & \text{judge class is \texttt{naive\_enough}},\\
1, & \text{judge class is \texttt{mix\_required}},
\end{cases}
\end{equation}
with \texttt{none\_enough} cases excluded from the binary router training set, since for these queries neither retrieval mode succeeds and the question is no longer a routing decision but an answer-quality failure.

The \gls{llmjudge} is treated as judged supervision rather than ground truth. The risk of judge error is mitigated in three ways. First, the judge sees the gold answer in addition to the candidate answers, which turns the task into a constrained comparison rather than open-ended quality scoring \cite{zheng2023judge}. Second, ambiguous \texttt{none\_enough} cases are excluded from router training rather than forced into one of the two binary classes. Third, the same judge model and prompt are used across the entire dataset, which keeps the labeling protocol internally consistent even if it is biased.

\section{Binary routing dataset}
\label{sec:methodDataset}

After the dual-mode execution and judge labeling, the binary routing dataset contains 1\,286 judged queries with a routing class of either \texttt{naive\_enough} or \texttt{mix\_required}. Of these, 959 (74.6\,\%) are labeled \texttt{naive\_enough} and 327 (25.4\,\%) are labeled \texttt{mix\_required}. The class distribution is therefore imbalanced, which is taken into account during training and during reporting (Section~\ref{sec:methodEvaluation}).

The dataset is split into 900 training examples, 193 validation examples, and 193 test examples (approximately 70/15/15), stratified by routing label so that both classes appear in all three splits in the same proportion as in the overall dataset. ``Held-out'' refers to data not used during training. The validation split is held out from training and is used only for early stopping and threshold selection. The test split is held out from both training and threshold selection and is used only for final reported metrics.

Table~\ref{tab:datasetSummary} summarizes the routing dataset by split and 2WikiMultihopQA question type. The class proportions are preserved across splits by construction, but the distribution of \texttt{mix\_required} labels is uneven across question types. \texttt{Compositional} questions are by far the most likely to require graph-enhanced retrieval, while \texttt{comparison} questions are usually answered by the cheap mode.

\begin{table}[!ht]
\begin{center}
\caption{Routing dataset summary by split and 2WikiMultihopQA question type. Counts show the number of judged binary examples in each cell; percentages in parentheses are the share of \texttt{mix\_required} within the cell.}
\label{tab:datasetSummary}
\begin{tabular}{l|r|rrrr}
\textbf{Split} & \textbf{Total} & \textbf{comparison} & \textbf{compositional} & \textbf{bridge\_comp.} & \textbf{inference} \\
\hline
Train & 900 & 411 (1.2\%) & 310 (58.4\%) & 153 (23.5\%) & 26 (26.9\%) \\
Validation & 193 & 94 (4.3\%) & 62 (64.5\%) & 29 (13.8\%) & 8 (12.5\%) \\
Test & 193 & 89 (3.4\%) & 64 (51.6\%) & 39 (33.3\%) & 1 (0.0\%) \\
\hline
Overall & 1\,286 & 594 (2.0\%) & 436 (58.3\%) & 221 (24.0\%) & 35 (22.9\%) \\
\end{tabular}
\end{center}
\end{table}

\section{Router models}
\label{sec:methodRouters}

Both router models take only the question text as input and output a score \(s_\phi(q_i) = P_\phi(y_i = 1 \mid q_i) \in [0,1]\) interpreted as the probability that \emph{Mix} is required.

\subsection{TF-IDF logistic-regression baseline}
\label{sec:methodTfidf}

The baseline router is a logistic regression trained on \gls{tfidf} features of the question. The pipeline first builds a vocabulary from the training questions, where each token is a unigram or bigram after lowercasing and basic English stop-word removal. Each question is then represented as a sparse vector \(\mathbf{x}_i \in \mathbb{R}^{V}\) whose entries are the term frequency of each vocabulary term multiplied by the inverse document frequency of that term in the training corpus. The logistic-regression head produces the router score as
\begin{equation}
\label{eq:tfidfScore}
s_\phi(q_i) = \sigma(\mathbf{w}^\top \mathbf{x}_i + b),
\qquad
\sigma(z) = \frac{1}{1 + e^{-z}},
\end{equation}
where \(\mathbf{w} \in \mathbb{R}^{V}\) and \(b \in \mathbb{R}\) are learned by minimizing the L2-regularized binary cross-entropy loss with balanced class weights, so that the minority \texttt{mix\_required} class is not ignored.

The baseline is intentionally simple. It has no pretrained semantic representations and can only exploit lexical patterns in the question. The role of this baseline is to indicate how much of the routing signal is captured by surface lexical cues, which calibrates how much improvement a transformer encoder can plausibly add.

\subsection{ModernBERT router}
\label{sec:methodModernBert}

The transformer router is a ModernBERT-base \cite{warner2024modernbert} encoder with a binary classification head. The question is tokenized and encoded by ModernBERT into a sequence of contextualized hidden states. The pooled representation \(h_i \in \mathbb{R}^{d}\) is then mapped to two logits by the classification head,
\begin{equation}
\label{eq:modernBertLogits}
z_i = W h_i + b,
\qquad
W \in \mathbb{R}^{2 \times d}, \quad b \in \mathbb{R}^{2},
\end{equation}
where \(z_{i,0}\) corresponds to \texttt{naive\_enough} and \(z_{i,1}\) to \texttt{mix\_required}. The router score is the softmax probability of the \texttt{mix\_required} class, given by
\begin{equation}
\label{eq:modernBertSoftmax}
s_\phi(q_i)
=
\frac{\exp(z_{i,1})}{\exp(z_{i,0}) + \exp(z_{i,1})}.
\end{equation}

The model is fine-tuned by minimizing class-balanced cross-entropy on the judged routing labels,
\begin{equation}
\label{eq:modernBertLoss}
\mathcal{L}(\phi)
=
-\frac{1}{n}\sum_{i=1}^{n}
\alpha_{y_i} \log P_\phi(y_i \mid q_i),
\end{equation}
where \(\alpha_{y_i}\) is the inverse class-frequency weight for the label of example \(i\), so that the minority \texttt{mix\_required} class is not down-weighted by the imbalanced training distribution. Optimization uses AdamW \cite{loshchilov2019decoupled} with a linear warm-up and linear decay schedule. The best checkpoint is selected on the validation split using the F1 score on the \texttt{mix\_required} class, and the chosen checkpoint is the one used for all reported test metrics. Hyperparameters are listed in Chapter~\ref{ch:implementation}.

\section{Evaluation}
\label{sec:methodEvaluation}

Evaluation has two complementary sides. Held-out classification metrics measure how well the router predicts the judged routing label. Offline policy evaluation measures what happens when the router's predictions are used to actually route queries between \emph{Naive} and \emph{Mix}. Both evaluations use the held-out test split. The validation split is used only for early stopping and for selecting an operating threshold.

\subsection{Held-out classification metrics}
\label{sec:methodHeldOut}

Because the binary class distribution is imbalanced, the primary threshold-independent metric is the \gls{auprc} on the \texttt{mix\_required} class. \Gls{auprc} measures how well the router ranks \texttt{mix\_required} queries above \texttt{naive\_enough} queries and is more informative than the area under the receiver-operating-characteristic curve when the positive class is in the minority \cite{saito2015precision}. The receiver-operating-characteristic area is also reported for comparison with prior work.

For thresholded evaluation, the router score \(s_\phi(q_i)\) is converted to a class prediction by applying a threshold \(\tau\). The thresholded policy is
\begin{equation}
\label{eq:routingPolicy}
\pi_\tau(q) =
\begin{cases}
M, & s_\phi(q) \geq \tau,\\
N, & s_\phi(q) < \tau,
\end{cases}
\end{equation}
where \(M\) denotes \emph{Mix} and \(N\) denotes \emph{Naive}. At a chosen threshold the reported classification metrics are accuracy, balanced accuracy, precision and recall on the \texttt{mix\_required} class, and the F1 score
\begin{equation}
\label{eq:f1}
F_1 = 2 \cdot
\frac{\mathrm{precision} \cdot \mathrm{recall}}
     {\mathrm{precision} + \mathrm{recall}}.
\end{equation}
Balanced accuracy is included because the class distribution is imbalanced and plain accuracy can be high even for a trivial always-\emph{Naive} policy.

Confidence intervals are reported using a 1\,000-iteration bootstrap that resamples the test rows with replacement and recomputes the metric on each resample. Reported intervals are the 2.5th and 97.5th percentiles of the bootstrap distribution. Multiple training seeds are also reported for the ModernBERT router, so that variance from random initialization can be separated from variance from finite test data.

\subsection{Offline routing policy}
\label{sec:methodOfflinePolicy}

The offline policy evaluation reuses the precomputed retrieval outputs and execution times from Section~\ref{sec:methodDualMode}. For each held-out query, both \(r_i^N\) and \(r_i^M\) are already available, along with their execution times \(t_i^N\) and \(t_i^M\). A routing policy \(\pi_\tau\) at threshold \(\tau\) is then simulated row by row. For each query the policy selects either the \emph{Naive} or the \emph{Mix} precomputed result, and aggregate statistics are computed across the entire held-out split. The mean routed execution time is
\begin{equation}
\label{eq:averageRoutedTime}
C(\pi_\tau) =
\frac{1}{m}\sum_{i=1}^{m}
\left[
\mathbf{1}\{\pi_\tau(q_i)=N\}\,t_i^N
+ \mathbf{1}\{\pi_\tau(q_i)=M\}\,t_i^M
\right],
\end{equation}
where \(m\) is the number of held-out queries. The same simulation is run for two static baseline policies for comparison. The first is an always-\emph{Naive} policy, and the second is an always-\emph{Mix} policy. These baselines bracket what a static retrieval strategy can achieve on the test set without any per-query decision and serve as the reference points against which the learned router's cost--quality trade-off is reported.

Threshold selection uses the validation split only. The threshold reported on test is the one that maximizes a chosen validation metric, and the same threshold is then applied unchanged to test. This prevents the threshold from being tuned on the data it is evaluated on.

\cleardoublepage
\chapter{Implementation}
\label{ch:implementation}

This chapter describes the implementation and the experiments. Section~\ref{sec:implTraining} describes how the two router models were trained. Section~\ref{sec:implClassification} presents the held-out classification experiment, which evaluates how well each router predicts the judged routing label. Section~\ref{sec:implPolicy} presents the cost--performance experiment, which evaluates each router as a thresholded routing policy against the static always-\emph{Naive} and always-\emph{Mix} baselines. Section~\ref{sec:implResults} reports the results of both experiments. Section~\ref{sec:implDiscussion} discusses what the results imply about the predictability of retrieval-mode selection, where the router fails, and what the practical cost--quality trade-off looks like.

\section{Model training}
\label{sec:implTraining}

Both routers are trained on the 900-example training split of the binary routing dataset described in Section~\ref{sec:methodDataset}. The validation split is used only for early stopping (ModernBERT) and threshold selection (both routers). The test split is used only for the metrics reported in Section~\ref{sec:implResults}.

\subsection{TF-IDF logistic-regression baseline}
\label{sec:implTfidf}

The baseline is built with scikit-learn. Questions are lowercased and stripped of standard English stop words, then converted into sparse term-frequency vectors with both unigram and bigram features and weighted by inverse document frequency learned on the training split. A logistic regression with $L_2$ regularization is fitted on top of these features, with the regularization strength left at the scikit-learn default and the class weights set to \texttt{balanced} so that the loss on each example is inversely proportional to the frequency of its class in the training data. Fitting on the 900 training questions takes a few seconds on a single CPU. The same fitted pipeline (vectorizer plus regression head) is then used to score the validation and test splits, producing a probability for the \texttt{mix\_required} class per query.

\subsection{ModernBERT router}
\label{sec:implModernBert}

The ModernBERT router uses the \texttt{answerdotai/ModernBERT-base} checkpoint from Hugging Face \cite{warner2024modernbert} with a binary classification head. Questions are tokenized with the standard ModernBERT tokenizer and truncated or padded to a maximum length of 128 tokens, which comfortably accommodates every question in the 2WikiMultihopQA splits used here.

Training uses the AdamW optimizer \cite{loshchilov2019decoupled} with a learning rate of $2 \times 10^{-5}$, a weight decay of $0.01$, gradient accumulation of two steps with a per-device batch size of $8$ (effective batch size $16$), a linear warm-up over the first $10\,\%$ of steps followed by linear decay, and gradient clipping at $1.0$. The loss is class-balanced cross-entropy with weights $\alpha_c$ inversely proportional to class frequency in the training split, so that the minority \texttt{mix\_required} class is not down-weighted by the imbalanced training distribution. The model is trained for up to $8$ epochs with early stopping based on validation F1 on the \texttt{mix\_required} class, requiring at least $3$ epochs before stopping and using a patience of $3$ epochs.

To separate variance from random initialization from variance from finite test data, the ModernBERT router is trained with three different random seeds ($7$, $13$, $42$). Each seed independently selects its own best checkpoint on the validation split and its own decision threshold on the validation split. The three runs are reported as a multi-seed mean and standard deviation in Section~\ref{sec:implResults}. All ModernBERT training and inference was performed on CPU. The model is small enough and the dataset short enough that the full fine-tuning run completes within a single afternoon, which is part of why a lightweight encoder router is attractive as a control layer.

\section{Held-out classification experiment}
\label{sec:implClassification}

The held-out classification experiment evaluates each router as a binary classifier of judged routing labels on the test split. The router scores are first evaluated threshold-independently, by computing \gls{auprc} on the \texttt{mix\_required} class and the area under the receiver-operating-characteristic curve (AUROC). \Gls{auprc} is the primary metric because it is sensitive to performance on the minority class, which is the class the routing decision actually cares about \cite{saito2015precision}.

For thresholded metrics, the decision threshold is selected on the validation split. The selection rule is the threshold $\tau^\star$ that maximizes F1 on \texttt{mix\_required} on the validation split, swept over a fixed grid in $[0.05, 0.95]$. The selected $\tau^\star$ is then applied unchanged to the test split, where accuracy, balanced accuracy, precision, recall, and F1 on \texttt{mix\_required} are reported. A trivial majority-class predictor (always-\emph{Naive}) is included as a degenerate reference because the class imbalance makes plain accuracy easy to inflate.

All metrics are reported with 95\,\% bootstrap intervals using 1\,000 resamples of the 193 test rows. For ModernBERT, point estimates are additionally reported across the three training seeds.

\section{Cost--performance experiment}
\label{sec:implPolicy}

The cost--performance experiment evaluates each router as an offline routing policy. The precomputed retrieval outputs and per-query execution times for both LightRAG modes are reused so that the simulation does not have to re-run any retrieval. A routing policy $\pi_{\tau}$ at threshold $\tau$ sends each test query either to the \emph{Naive} or to the \emph{Mix} precomputed result. The mean routed time on the test split is then computed as in Equation~\ref{eq:averageRoutedTime}.

For each router, the threshold curve is computed over a dense grid of $\tau$ values in $[0, 1]$, producing a Pareto-style trace in (mean routed time, F1\,on \texttt{mix\_required}) space. The same simulation is run for the two static baseline policies, always-\emph{Naive} and always-\emph{Mix}. These bracket what a static policy can achieve on the test set and serve as the reference points against which the learned routers are compared.

\section{Results}
\label{sec:implResults}

Table~\ref{tab:headlineResults} summarizes the test-set performance of the four reported policies. The first two are the static always-\emph{Naive} and always-\emph{Mix} baselines. The third is the TF-IDF baseline router at its validation-selected threshold. The fourth is the ModernBERT router at its validation-selected threshold, averaged over three seeds. Bootstrap confidence intervals are shown for the learned routers' key metrics. The ModernBERT router improves on the TF-IDF baseline on every metric except AUROC, where the difference is within the bootstrap interval. The largest gain is in balanced accuracy (from $0.727$ to $0.756$) and recall on \texttt{mix\_required} (from $0.531$ to $0.605$), which together indicate that the encoder router catches more of the queries that genuinely need graph-enhanced retrieval. Both learned routers leave the always-\emph{Mix} policy far behind on cost without losing much accuracy.

\begin{table}[!ht]
\centering
\caption{Test-set performance summary. \texttt{AUPRC} and \texttt{AUROC} are threshold-independent. \texttt{F1\textsubscript{mix}}, balanced accuracy, accuracy, mean routed time and mix-route rate are reported at the validation-selected threshold $\tau^\star$. Bracketed values for the learned routers are 95\,\% bootstrap intervals on the test set (1\,000 resamples). ModernBERT values are the mean over three training seeds $\pm$ standard deviation.}
\label{tab:headlineResults}
\resizebox{\textwidth}{!}{%
\begin{tabular}{l|cccccc}
\textbf{Policy} & \textbf{AUPRC} & \textbf{AUROC} & \textbf{Bal.\ acc.} & \textbf{F1\textsubscript{mix}} & \textbf{Acc.} & \textbf{Mean time (s)} \\
\hline
Always-\emph{Naive} & --- & --- & 0.500 & 0.000 & 0.746 & 10.88 \\
Always-\emph{Mix}   & --- & --- & 0.500 & 0.405 & 0.254 & 21.36 \\
TF-IDF (LR) & 0.708 [0.581, 0.821] & 0.864 [0.798, 0.917] & 0.727 & 0.605 [0.470, 0.722] & 0.824 [0.768, 0.876] & 13.03 \\
ModernBERT (mean $\pm$ s.d.) & $0.732 \pm 0.005$ & $0.854 \pm 0.004$ & $0.756 \pm 0.008$ & $0.645 \pm 0.002$ & $0.831 \pm 0.013$ & $13.94 \pm 0.17$ \\
\end{tabular}%
}
\end{table}

Figure~\ref{fig:prCurveTest} shows the test-set precision--recall curves on the \texttt{mix\_required} class for the TF-IDF baseline and the three ModernBERT seeds. The base rate (the test-set fraction of \texttt{mix\_required}) is shown as a horizontal dashed line. The ModernBERT curves dominate the TF-IDF curve over most of the recall range, which is consistent with the AUPRC numbers in Table~\ref{tab:headlineResults}. The receiver-operating-characteristic curves (Figure~\ref{fig:rocCurveTest}) are closer together, reflecting that AUROC is less sensitive to performance on the minority class.

\begin{figure}[!ht]
\centering
\includegraphics[width=0.85\textwidth]{figures/thesis/pr_curve_test.pdf}
\caption{Precision--recall curves on the held-out test split for the TF-IDF baseline and the three ModernBERT training seeds. The dashed horizontal line is the test-set base rate of \texttt{mix\_required}.}
\label{fig:prCurveTest}
\end{figure}

\begin{figure}[!ht]
\centering
\includegraphics[width=0.85\textwidth]{figures/thesis/roc_curve_test.pdf}
\caption{Receiver-operating-characteristic curves on the held-out test split.}
\label{fig:rocCurveTest}
\end{figure}

Figure~\ref{fig:thresholdVsF1} shows how F1 on \texttt{mix\_required} varies with the decision threshold on the test split, with the validation-selected threshold of each model marked as a dot. The ModernBERT seeds disagree about where their optimal validation threshold lies (seed $42$ selects $\tau^\star = 0.24$, seed $13$ selects $\tau^\star = 0.53$, seed $7$ selects $\tau^\star = 0.83$), but their test F1 at the chosen threshold is consistent (mean $0.645$, standard deviation $0.002$), suggesting the F1 plateau is broad and the routers degrade gracefully when the threshold is mistuned.

\begin{figure}[!ht]
\centering
\includegraphics[width=0.85\textwidth]{figures/thesis/threshold_vs_f1_mix.pdf}
\caption{F1 on \texttt{mix\_required} as a function of the decision threshold on the held-out test split. Dots mark the validation-selected threshold for each model.}
\label{fig:thresholdVsF1}
\end{figure}

Figure~\ref{fig:thresholdVsTime} shows the mean routed time on test as the threshold sweeps from $0$ (every query routed to \emph{Mix}) to $1$ (every query routed to \emph{Naive}). The always-\emph{Naive} and always-\emph{Mix} levels are shown as horizontal dashed lines and act as a sanity check. At $\tau = 0$ every router collapses to always-\emph{Mix}, and at $\tau = 1$ every router collapses to always-\emph{Naive}, which is what the figure shows.

\begin{figure}[!ht]
\centering
\includegraphics[width=0.85\textwidth]{figures/thesis/threshold_vs_routed_time.pdf}
\caption{Mean routed execution time on the held-out test split as the decision threshold sweeps. The two horizontal dashed lines are the always-\emph{Naive} and always-\emph{Mix} reference levels.}
\label{fig:thresholdVsTime}
\end{figure}

Figure~\ref{fig:costFrontier} combines the two preceding views into a cost--performance frontier. Each router traces a Pareto curve in (mean routed time, F1\textsubscript{mix}) space as $\tau$ sweeps. The static always-\emph{Naive} and always-\emph{Mix} policies appear as the two extreme operating points. The validation-selected operating points for the learned routers lie above the straight line connecting the two static endpoints. This is the operating range in which a router is genuinely useful, providing operating points that no static policy can reach by itself. The ModernBERT router lies slightly above and to the right of the TF-IDF baseline, reaching higher F1\textsubscript{mix} at a slightly higher mean routed time.

\begin{figure}[!ht]
\centering
\includegraphics[width=0.85\textwidth]{figures/thesis/cost_performance_frontier.pdf}
\caption{Cost--performance frontier on the held-out test split. Each curve traces F1 on \texttt{mix\_required} versus mean routed execution time as the decision threshold sweeps. Stars mark the validation-selected operating point of each router. The black square (always-\emph{Naive}) and the black diamond (always-\emph{Mix}) are the two static baselines.}
\label{fig:costFrontier}
\end{figure}

Table~\ref{tab:perQuestionType} reports a per-question-type breakdown on the test split. The 2WikiMultihopQA question types have very different routing distributions. The \texttt{comparison} slice is almost always \texttt{naive\_enough}, with only $3$ of $89$ test rows requiring \emph{Mix}. The \texttt{bridge\_comparison} slice sits around one third. The \texttt{compositional} slice is the one in which the routing decision matters most, with $33$ of $64$ test rows requiring \emph{Mix}. The \texttt{inference} slice has only one test row and is omitted from the per-class metrics. On the \texttt{compositional} slice the ModernBERT router improves AUPRC over TF-IDF (mean $0.846$ vs. $0.787$). On \texttt{bridge\_comparison} the two routers behave qualitatively differently. At the validation-selected threshold the TF-IDF router never predicts \emph{Mix} on this slice (recall $0$), while ModernBERT recovers some examples (mean recall $0.154$, mean F1 $0.245$). The AUPRC, however, is slightly lower for ModernBERT, which means TF-IDF ranks the few positive bridge\_comparison rows above its threshold but never crosses the threshold itself.

\begin{table}[!ht]
\centering
\caption{Per-question-type breakdown on the test split. The \texttt{comparison} and \texttt{inference} slices have $\leq 3$ positives and are reported only for the share of \texttt{mix\_required}. ModernBERT values are the mean over three seeds.}
\label{tab:perQuestionType}
\begin{tabular}{l|rrrrrr}
\textbf{Slice} & \textbf{n} & \textbf{n\textsubscript{mix}} & \textbf{Model} & \textbf{AUPRC} & \textbf{F1\textsubscript{mix}} & \textbf{Acc.} \\
\hline
comparison        &  89 &  3 & TF-IDF     & ---    & 0.000 & 0.966 \\
                  &     &    & ModernBERT & ---    & 0.000 & 0.966 \\
\hline
compositional     &  64 & 33 & TF-IDF     & 0.787 & 0.743 & 0.719 \\
                  &     &    & ModernBERT & 0.846 & 0.762 & 0.729 \\
\hline
bridge\_comparison&  39 & 13 & TF-IDF     & 0.646 & 0.000 & 0.667 \\
                  &     &    & ModernBERT & 0.517 & 0.245 & 0.684 \\
\hline
inference         &   1 &  0 & ---        & ---   & ---   & 1.000 \\
\end{tabular}
\end{table}

Figure~\ref{fig:perQt} visualizes the AUPRC comparison on the two non-degenerate slices.

\begin{figure}[!ht]
\centering
\includegraphics[width=0.85\textwidth]{figures/thesis/per_question_type_auprc.pdf}
\caption{AUPRC on the \texttt{mix\_required} class, broken down by 2WikiMultihopQA question type on the held-out test split. Error bars on ModernBERT are the standard deviation across three training seeds. The \texttt{comparison} and \texttt{inference} slices have too few \texttt{mix\_required} test rows to give an informative AUPRC and are omitted.}
\label{fig:perQt}
\end{figure}

\section{Discussion}
\label{sec:implDiscussion}

The headline result is that the retrieval-mode distinction between \emph{Naive} and \emph{Mix} is partly predictable from the query text alone. Both routers comfortably outperform the always-\emph{Naive} majority baseline on the minority class and reach operating points that are not on the straight line between the two static policies in Figure~\ref{fig:costFrontier}. The ModernBERT router is consistently better than the TF-IDF baseline on AUPRC, on balanced accuracy, and on recall of the minority class, with differences that are stable across three training seeds.

The most informative slice is \texttt{compositional}. This is where the class distribution is genuinely balanced and where graph-enhanced retrieval is most often needed. On this slice ModernBERT achieves a mean AUPRC of $0.846$ against TF-IDF's $0.787$, despite the small training set, which suggests that the encoder picks up on more than the lexical patterns the bag-of-words classifier can exploit. The \texttt{bridge\_comparison} slice is harder. TF-IDF achieves higher AUPRC but never crosses its own validation-selected threshold, so it routes everything in this category to \emph{Naive} and the \texttt{mix\_required} recall is exactly zero. ModernBERT scores fewer queries in this slice as confidently positive but is at least willing to cross its threshold, which yields a non-trivial F1 there. Both behaviors are honest reflections of how the routers handle that slice, and neither is good. Improving routing on this slice is an obvious target for further work.

There are limitations to the headline result. The held-out test set has 193 queries, which is enough to compute meaningful bootstrap intervals but not enough to resolve small differences between the routers. The AUROC intervals of TF-IDF ($[0.798, 0.917]$) and ModernBERT ($\approx [0.78, 0.92]$ for each seed) overlap, so the AUROC comparison cannot be called significant. The AUPRC and balanced-accuracy gaps are larger but still inside the bootstrap intervals, which is what should be expected on a $193$-row test set. The multi-seed standard deviation for ModernBERT is small ($0.005$ on AUPRC, $0.008$ on balanced accuracy), so the gains are at least stable across seeds, but a larger benchmark would be needed before stronger claims can be made.

The cost picture is more clear-cut. The two static baselines bracket a mean routed time range of roughly $10.9$ to $21.4$ seconds per query. The learned routers select operating points around $13$ to $14$ seconds with non-trivial \texttt{mix\_required} F1, which corresponds to roughly $35$ to $40\,\%$ less mean routed time than always-\emph{Mix} at a much higher F1 on the minority class than always-\emph{Naive}. The routing decision itself is essentially free. A single forward pass through a $\sim$100\,M-parameter encoder, or a sparse logistic-regression evaluation, is negligible next to the seconds-scale retrieval cost it replaces. Section~\ref{sec:methodEvaluation} discussed why execution time was chosen as the cost metric rather than token usage; reporting a token-level cost would have required instrumenting LightRAG's internal language-model calls, which was outside the scope of this thesis.

\cleardoublepage
\chapter{Conclusions and future work}
\label{ch:conclusionsAndFutureWork}

This chapter summarizes the conclusions from the experiments (Section~\ref{sec:conclusions}), discusses the main limitations and natural avenues for future work (Section~\ref{sec:limitationsFutureWork}), and reflects on the broader context of the project (Section~\ref{sec:reflections}).

\section{Conclusions}
\label{sec:conclusions}

This thesis studied whether the choice between LightRAG's cheap \emph{Naive} retrieval mode and its more expensive \emph{Mix} retrieval mode can be predicted from the query text alone, on the 2WikiMultihopQA benchmark, using a lightweight encoder-only classifier rather than a large language-model controller. Routing labels were derived by executing both retrieval modes on each benchmark question and using an \gls{llmjudge} to compare each mode's answer against the gold answer.

On the held-out test split, the ModernBERT router improves on a TF-IDF logistic-regression baseline on \gls{auprc} ($0.732$ vs. $0.708$), balanced accuracy ($0.756$ vs. $0.727$), F1 on the minority \texttt{mix\_required} class ($0.645$ vs. $0.605$), and recall on the minority class ($0.605$ vs. $0.531$), while remaining within the bootstrap interval on AUROC. The multi-seed standard deviation of the ModernBERT router is small ($\leq 0.013$ on every reported metric), which makes the improvement stable rather than seed-specific. Both learned routers reach test-set operating points around $13$--$14$ seconds of mean routed execution time, compared with $10.9$ seconds for always-\emph{Naive} and $21.4$ seconds for always-\emph{Mix}, and lie above the straight line connecting the two static endpoints in the cost--performance frontier.

The per-question-type breakdown localizes most of the routing signal in the \texttt{compositional} slice, where graph-enhanced retrieval is most often needed (51\,\% of the test slice) and where ModernBERT improves AUPRC over TF-IDF by roughly six points. The \texttt{comparison} slice is essentially decidable as \emph{Naive} and both routers correctly avoid \emph{Mix} on it. The \texttt{bridge\_comparison} slice is the weakest part of the result and is the natural target for further work.

Returning to the research question, the distinction between vector-sufficient and graph-required queries is, in this setting, partly predictable from the query text alone. The gain over a classical baseline is small but stable, and the practical effect is non-trivial cost savings at a level of \texttt{mix\_required} recall that a static policy cannot match.

\section{Limitations and future work}
\label{sec:limitationsFutureWork}

\subsection{Token-level cost}
\label{sec:limitTokenCost}

The cost metric reported throughout the thesis is end-to-end execution time per query. This was the cleanly measurable cost because both LightRAG modes were already timed at the per-query level. Token usage would be a natural alternative cost dimension and is what users pay for in production language-model APIs, but LightRAG's modes invoke an internal language model several times during retrieval (for graph traversal, context assembly, and entity expansion) and these internal calls were not instrumented in this work. A natural follow-up is to instrument LightRAG's internal calls and report both token cost and execution time as parallel cost dimensions, which would give a clearer picture of how much of the saving comes from generation-side tokens versus retrieval-side language-model calls.

\subsection{Judge labels}
\label{sec:limitJudgeLabels}

Routing labels were produced by a single \gls{llmjudge} model with access to the gold answer. This is a reasonable but imperfect labeling protocol \cite{zheng2023judge}. The risk is that the judge systematically misclassifies particular kinds of correct answers, which would shift the routing labels in a way that the router learns to mimic. A natural follow-up is to label a sample of queries with a different judge model or with human annotators and measure inter-rater agreement on the routing label, both as a sanity check on the existing labels and as a more reliable evaluation subset.

\subsection{Generalizability}
\label{sec:limitGeneralizability}

The experiments use one Hybrid \gls{rag} framework (LightRAG), two of its retrieval modes (\emph{Naive} and \emph{Mix}), and one public benchmark (2WikiMultihopQA). The held-out test set has $193$ rows, which is enough to compute bootstrap intervals but not enough to resolve small differences between routers. The conclusions therefore apply most directly to the studied setting, and applying the same methodology to other Hybrid \gls{rag} frameworks (for example GraphRAG or Hybrid retrieval modes in other systems) or to other multi-hop benchmarks (HotpotQA, MuSiQue) is a natural way to test how robust the routing signal is. Extending the routing labels to industrial corpora (for example Ericsson internal documentation) was considered but is blocked by the absence of per-query ground-truth answers in those corpora, which is what made the judge-labeling protocol possible on 2WikiMultihopQA in the first place.

\subsection{Routing beyond the query}
\label{sec:limitRouting}

The router used in this thesis is query-only by construction. A router that also has access to some lightweight signal from the corpus or from a cheap first retrieval pass could in principle make a better decision, at the cost of a slightly more expensive control layer. A natural follow-up is to compare the query-only router against a hybrid scheme that first runs the cheap \emph{Naive} mode, inspects whether the retrieved chunks contain the entities referenced in the query, and then escalates to \emph{Mix} only when this signal looks weak. This would still keep the control layer cheap (one round of vector retrieval rather than an expensive graph traversal) but could materially improve the harder slices, in particular \texttt{bridge\_comparison}.

\subsection{Calibration and abstention}
\label{sec:limitCalibration}

The router currently outputs a single thresholded decision. Routing under cost constraints is a natural setting for selective prediction \cite{geifman2017selective} and calibrated probabilities \cite{guo2017calibration}, where the router could abstain on low-confidence queries (for example by escalating to a more expensive but more reliable controller, or by running both retrieval modes and combining the answers). The calibration of the ModernBERT router's softmax outputs was not investigated in this work and is a natural piece of follow-up.

\section{Reflections}
\label{sec:reflections}

This thesis is centered on making a Hybrid \gls{rag} system more resource-efficient. The most direct connection is to the \gls{UN} \gls{SDG} number 12 (responsible consumption and production). Graph-enhanced retrieval is significantly more compute-intensive than vector retrieval, and a router that avoids applying it to queries that do not need it directly reduces the energy and compute footprint of the system at inference time. There is also a softer connection to \gls{SDG} number 9 (industry, innovation, and infrastructure), in that adaptive routing is a building block for industrial \gls{rag} deployments such as the one motivating this work at Ericsson, where serving company-internal documentation at scale is a real operational concern.

The work has some ethical aspects worth flagging. The routing labels are produced by an \gls{llmjudge} that has access to the gold answer. This is a methodologically reasonable but imperfect protocol, and the thesis is explicit about not treating these labels as ground truth. Building a system that decides on a per-query basis which retrieval path to use could in principle interact with downstream notions of answer quality in ways that are hard to anticipate. If the router systematically routes some kind of question to the cheaper path, that kind of question could end up consistently worse-served, even when the average behavior of the system improves. The per-question-type breakdown in Section~\ref{sec:implResults} is one way to surface this kind of differential behavior, and any future work that scales this approach to user-facing settings should keep similar per-segment analyses in view.

Finally, the implementation code, the data-preparation pipeline, the training scripts, and the evaluation scripts used in this thesis are available at \url{\thesisCodeUrl}.

\cleardoublepage
% Print the bibliography (and make it appear in the table of contents)
\renewcommand{\bibname}{References}


\iftoggle{biblatex}{
    %\typeout{Biblatex current language is \currentlang}
    \printbibliography[heading=bibintoc]
}{
    \phantomsection  % make it include a hyperref - see https://tex.stackexchange.com/a/98995
    \addcontentsline{toc}{chapter}{References}
    \bibliography{references}
}



\cleardoublepage
\appendix
\renewcommand{\chaptermark}[1]{\markboth{Appendix \thechapter\relax:\thinspace\relax#1}{}}
\chapter{Supporting materials}
\label{sec:supportingMaterial}

The BibTeX references used in this thesis are attached to this PDF.
\attachfile[description={references.bib}]{references.bib}

The implementation code, the dataset-preparation scripts, the LightRAG indexing and querying pipeline, the LLM-as-a-Judge labeling pipeline, the router training scripts, and the analysis and plotting scripts used in this thesis are available at \url{\thesisCodeUrl}.

%% The following label is necessary for computing the last page number of the body of the report to include in the "For DIVA" information
\label{pg:lastPageofMainmatter}

\cleardoublepage
\clearpage\thispagestyle{empty}\mbox{} % empty page with backcover on the other side
\kthbackcover
\fancyhead{}  % Do not use header on this extra page or pages
\section*{€€€€ For DIVA €€€€}
\lstset{numbers=none} %% remove any list line numbering
\divainfo{pg:lastPageofPreface}{pg:lastPageofMainmatter}

% If there is an acronyms.tex file,
% add it to the end of the For DIVA information
% so that it can be used with the abstracts
% Note that the option "nolol" stops it from being listed in the List of Listings

% The following bit of ugliness is because of the problems PDFLaTeX has handling a non-breaking hyphen
% unless it is converted to UTF-8 encoding.
% If you do not use such characters in your acronyms, this could be simplified.
\ifxeorlua
\IfFileExists{lib/acronyms.tex}{
\section*{acronyms.tex}
\lstinputlisting[language={[LaTeX]TeX}, nolol, basicstyle=\ttfamily\color{black},
commentstyle=\color{black}, backgroundcolor=\color{white}]{lib/acronyms.tex}
}
{}
\else
\IfFileExists{lib/acronyms-for-pdflatex.tex}{
\section*{acronyms.tex}
\lstinputlisting[language={[LaTeX]TeX}, nolol, basicstyle=\ttfamily\color{black},
commentstyle=\color{black}, backgroundcolor=\color{white}]{lib/acronyms-for-pdflatex.tex}
}
{}
\fi


\end{document}
